% =========================================================================
% Atelier IA — Lead-magnet PDF
% « Claude Desktop au quotidien »
%
% Compile : pdflatex -output-directory=build lead_magnet.tex (deux passes)
% =========================================================================

\documentclass[11pt,a4paper]{article}

% --- Encoding & language --------------------------------------------------
\usepackage[T1]{fontenc}
\usepackage[utf8]{inputenc}
\usepackage[french]{babel}
\frenchsetup{ItemLabeli=\textendash, ItemLabelii=\textendash}

% --- Typography -----------------------------------------------------------
\usepackage{ebgaramond}              % corps : Garamond ouvert et lisible
\usepackage[scale=0.85]{inconsolata} % mono pour datelines / labels
\usepackage{microtype}               % protrusion + expansion typo

% Pour les titres : reprend le family Cormorant via fontspec ? Non — on reste
% en EB Garamond italique pour les italiques de titre (cohérent avec le rendu
% du site qui utilise Instrument Serif italic).

% --- Geometry -------------------------------------------------------------
\usepackage[
  a4paper,
  top=3cm, bottom=3cm,
  left=3cm, right=3cm,
  headsep=1cm, footskip=1.5cm
]{geometry}

% --- Colors (matchent la landing : vélin, encre, bordeaux) ---------------
\usepackage{xcolor}
\definecolor{velin}{HTML}{F5EFE2}
\definecolor{encre}{HTML}{1A1614}
\definecolor{encredoux}{HTML}{3D362F}
\definecolor{encrefaint}{HTML}{8A8077}
\definecolor{bordeaux}{HTML}{7A2E2A}

% Application : texte body en encre douce
\AtBeginDocument{\color{encredoux}}

% --- Section formatting ---------------------------------------------------
\usepackage{titlesec}

% Suppress section numbering (we use our own roman numerals editorially)
\setcounter{secnumdepth}{0}

\titleformat{\section}[block]
  {\color{encre}\normalfont\Huge\itshape}
  {}{0em}{}
\titlespacing*{\section}{0pt}{2em}{1em}

\titleformat{\subsection}[block]
  {\color{encre}\normalfont\Large\bfseries}
  {}{0em}{}
\titlespacing*{\subsection}{0pt}{1.5em}{0.5em}

\titleformat{\subsubsection}[block]
  {\color{bordeaux}\normalfont\small\sffamily\MakeUppercase}
  {}{0em}{}
\titlespacing*{\subsubsection}{0pt}{1em}{0.3em}

% --- Lists ---------------------------------------------------------------
\usepackage{enumitem}
\setlist{leftmargin=1.5em, itemsep=0.3em, parsep=0pt, topsep=0.4em}

% --- Headers / footers ---------------------------------------------------
\usepackage{fancyhdr}
\pagestyle{fancy}
\fancyhf{}
\renewcommand{\headrulewidth}{0pt}
\renewcommand{\footrulewidth}{0.4pt}
\fancyfoot[L]{\footnotesize\ttfamily\color{encrefaint}Atelier IA}
\fancyfoot[C]{\footnotesize\ttfamily\color{encrefaint}Claude Desktop au quotidien}
\fancyfoot[R]{\footnotesize\ttfamily\color{encrefaint}\thepage}

% Cover page : no header/footer
\fancypagestyle{cover}{
  \fancyhf{}
  \renewcommand{\headrulewidth}{0pt}
  \renewcommand{\footrulewidth}{0pt}
}

% --- Custom commands -----------------------------------------------------
\newcommand{\dateline}[1]{{\footnotesize\ttfamily\color{encrefaint}\MakeUppercase{#1}}}
\newcommand{\bordot}{\textcolor{bordeaux}{$\bullet$}}
\newcommand{\bordsec}[1]{\subsubsection*{\textcolor{bordeaux}{#1}}}

% Pull-quote / aside
\newcommand{\pullquote}[1]{%
  \begin{quote}\itshape\color{encre}\large #1\end{quote}%
}

% Skill block — formatted like a code listing
\usepackage{tcolorbox}
\tcbuselibrary{skins}
\newtcolorbox{skillbox}[1]{
  enhanced,
  colback=velin,
  colframe=encre,
  boxrule=0.4pt,
  arc=0pt,
  outer arc=0pt,
  fonttitle=\ttfamily\small\color{encre},
  title={\MakeUppercase{Skill — #1}},
  attach boxed title to top left={xshift=0pt, yshift=-3pt},
  boxed title style={colback=velin, colframe=velin, boxrule=0pt, arc=0pt},
  left=1em, right=1em, top=1.2em, bottom=1em,
  fontupper=\small\color{encredoux}
}

% Workspace block
\newtcolorbox{workspacebox}{
  enhanced,
  colback=velin,
  colframe=bordeaux,
  boxrule=0pt,
  leftrule=2pt,
  arc=0pt,
  outer arc=0pt,
  left=1em, right=0.5em, top=0.8em, bottom=0.8em,
  fontupper=\small\color{encredoux}
}

% Hyperlinks
\usepackage[hidelinks,pdfusetitle]{hyperref}
\hypersetup{
  pdfauthor={Vincent Tariel},
  pdftitle={Atelier IA — Claude Desktop au quotidien},
  pdfsubject={Guide pratique pour intégrer Claude Desktop dans son métier},
  pdfkeywords={Claude, Desktop, Coaching, IA, Nouméa, Calédonie}
}

% Tighten paragraphs
\setlength{\parskip}{0.5em}
\setlength{\parindent}{0pt}

% =========================================================================
\begin{document}
% =========================================================================

% -------------------------------------------------------------------------
% COVER
% -------------------------------------------------------------------------
\thispagestyle{cover}

\vspace*{2cm}

\dateline{Atelier IA \quad\textbar\quad Numéro 01 \quad\textbar\quad Édition 2026}

\vspace{0.3cm}
{\color{encre}\rule{\linewidth}{1pt}}
\vspace{1.5cm}

\begin{flushleft}
{\color{encre}\fontsize{52pt}{56pt}\selectfont\itshape Claude Desktop\par}
\vspace{0.4cm}
{\color{encre}\fontsize{52pt}{56pt}\selectfont au quotidien\textcolor{bordeaux}{.}\par}
\end{flushleft}

\vspace{1.5cm}

\begin{minipage}{0.65\textwidth}
{\large\color{encredoux}
Le guide pratique pour les professionnels qui veulent
faire travailler l'\textit{intelligence artificielle d'Anthropic}
sur leurs vrais dossiers — sans coder.
}
\end{minipage}

\vfill

{\color{encre}\rule{\linewidth}{0.4pt}}

\vspace{0.3cm}

\begin{minipage}{0.6\textwidth}
\dateline{Vincent Tariel}\\[0.2em]
{\footnotesize\itshape Docteur en mathématiques de l'École Polytechnique\\
Coaching Claude Desktop, en présentiel à Nouméa\\
\texttt{atelier-ia.ovh}}
\end{minipage}\hfill
\begin{minipage}{0.35\textwidth}\raggedleft
\dateline{À lire en 12 minutes}\\[0.2em]
{\footnotesize\itshape Premier guide d'une série en cours.}
\end{minipage}

\newpage

% -------------------------------------------------------------------------
% I. PRÉFACE
% -------------------------------------------------------------------------

\dateline{I. \quad Préface}

\vspace{0.4cm}
{\color{encre}\rule{\linewidth}{0.4pt}}
\vspace{1cm}

\section*{Pourquoi ce guide existe}

Vous savez probablement déjà utiliser une intelligence artificielle. Vous avez
ouvert ChatGPT ou Claude dans un onglet de votre navigateur, posé vos questions,
obtenu des réponses utiles. Le chat web fait très bien son travail — pour des
questions ponctuelles.

Le problème commence quand vous voulez l'intégrer \textit{sérieusement} à
votre quotidien. À chaque nouvelle session, il faut redéposer vos fichiers.
Réécrire le contexte. Retaper le rôle attendu. Cinq à dix minutes de mise en
route pour deux minutes de réponse utile. Multipliée par les jours, cette
répétition tue l'usage — et vous finissez par retourner à Word, à votre boîte
mail, à votre traitement de texte habituel.

\textit{Pourquoi ne pas travailler directement sur votre ordinateur~?} Avec
vos fichiers à leur vraie place, votre arborescence, vos modèles
immédiatement accessibles — sans rien à uploader à chaque ouverture, sans
rien à ré-expliquer~?

C'est exactement ce que \textbf{Claude Desktop} apporte. L'application,
installée sur votre Mac ou votre PC, accède à vos fichiers locaux. Une fois
configurée, elle se souvient de votre rôle, de vos références, de vos
contraintes. Vous l'ouvrez, vous lancez la tâche, vous récupérez un
brouillon prêt à corriger. Le chat web devient un outil de dépannage~;
Claude Desktop devient \textit{l'outil de travail}.

\textbf{Ce guide n'est pas un manuel.} C'est un récit en deux histoires —
celle d'un professeur de mathématiques au lycée, celle d'un avocat débordé
par sa rédaction. Vous y verrez comment chacun~:

\begin{itemize}
  \item construit son \textit{espace de travail} dans Claude Desktop,
  \item définit un \textit{skill} appelable à la demande,
  \item s'épargne, semaine après semaine, des heures de travail sans valeur
        ajoutée.
\end{itemize}

À la fin du guide, vous aurez compris la méthode. Vous saurez si elle peut
s'appliquer à votre métier. Si oui, mes coordonnées sont en dernière page.

\pullquote{« Claude ne remplace pas votre travail. Il augmente votre productivité
pour que vous restiez focus sur votre vraie valeur ajoutée. »}

\newpage

% -------------------------------------------------------------------------
% II. LA MÉTHODE
% -------------------------------------------------------------------------

\dateline{II. \quad La méthode}

\vspace{0.4cm}
{\color{encre}\rule{\linewidth}{0.4pt}}
\vspace{1cm}

\section*{Deux temps, et c'est tout}

Toute la méthode tient en deux mouvements. C'est rare dans la sphère IA, où
les guides s'épaississent à mesure qu'ils s'éloignent du métier de leurs
lecteurs. Ici, deux pages suffiront.

\bordsec{Temps 1 — L'espace de travail}

Claude Desktop fonctionne avec ce qu'on appelle des \textit{Projects}.
Un Project, c'est un dossier dans lequel vous déposez les documents de
référence d'un domaine précis~: vos modèles de courriers, vos textes officiels,
vos derniers rapports, votre glossaire interne. Une fois ces documents
en place, Claude les a en permanence sous les yeux quand vous lui parlez
dans ce Project.

Le travail consiste à choisir \textit{quoi} mettre dans cet espace.
C'est moins évident qu'il n'y paraît. Mettre trop de documents noie le
contexte~; trop peu produit des réponses creuses. Le bon dosage demande
une heure de travail au début — et tient ensuite des mois.

\bordsec{Temps 2 — Le skill appelable}

Une fois l'espace de travail prêt, on définit un \textit{skill}~: une
consigne précise et permanente qui décrit \textbf{le rôle} que Claude doit
tenir dans ce Project. Le skill répond à trois questions~:

\begin{itemize}
  \item Qui es-tu~? (\textit{« Tu es l'assistant de rédaction d'un cabinet d'avocats à Nouméa »})
  \item Comment travailles-tu~? (\textit{« Tu reproduis le style des modèles fournis, tu numérotes les pièces, tu poses une question si une information manque »})
  \item Qu'est-ce que tu refuses de faire~? (\textit{« Tu n'inventes jamais de jurisprudence sans source explicite »})
\end{itemize}

Ce skill, une fois enregistré, devient appelable à toute heure du jour
sans rien réexpliquer. Vous lancez le Project, vous donnez votre tâche
brute (\textit{« voici les notes de mon RDV de ce matin, rédige
l'attestation »}), et vous récupérez un brouillon en trente secondes.

\vspace{1cm}

\pullquote{Espace de travail bien rempli, skill bien rédigé, patience
au début~: ce sont les trois investissements qui transforment Claude
en collaborateur — et non en jouet.}

\newpage

% -------------------------------------------------------------------------
% III. HISTOIRE 1 — LE PROFESSEUR
% -------------------------------------------------------------------------

\dateline{III. \quad Histoire 1}

\vspace{0.4cm}
{\color{encre}\rule{\linewidth}{0.4pt}}
\vspace{1cm}

\section*{Le professeur de mathématiques}

\subsubsection*{Le contexte}

Marc enseigne les mathématiques en seconde générale au lycée Jules Garnier
à Nouméa. Trente-deux élèves par classe, deux classes en parallèle, un
programme officiel chargé. Chaque dimanche soir, il prépare ses cours pour
la semaine~: structure, exercices, évaluations formatives, différenciation
pour les élèves en difficulté. Trois heures de préparation pour cinq heures
de cours hebdomadaires. Multiplié par les semaines, c'est une nuit complète
qu'il aimerait récupérer.

\bordsec{Temps 1 — Son espace de travail}

Marc crée un Project Claude Desktop nommé \textit{« Maths · Seconde
générale »}. Il y dépose~:

\begin{workspacebox}
\begin{itemize}[leftmargin=1.2em]
  \item Le \textbf{Bulletin Officiel} 2024 du programme de seconde générale
        (PDF récupéré sur le site du Ministère).
  \item Les \textbf{progressions annuelles} de deux collègues plus expérimentés
        — récupérées en salle des profs, avec leur accord.
  \item Trois \textbf{manuels scolaires} de référence (versions PDF
        légalement obtenues via le portail enseignant).
  \item Ses \textbf{cours de l'an dernier}, dans leur état brut.
  \item Un fichier \textbf{« Mes 32 élèves »} anonymisé~: prénoms remplacés
        par des codes, niveau estimé, points à surveiller.
\end{itemize}
\end{workspacebox}

Ce travail de mise en place lui prend une heure et demie un samedi matin.
Il ne le refera plus de l'année.

\bordsec{Temps 2 — Son skill}

Marc rédige ensuite l'instruction permanente du Project~:

\begin{skillbox}{Professeur de mathématiques}
Tu es professeur de mathématiques en seconde générale au lycée Jules Garnier
à Nouméa. Tu maîtrises le Bulletin Officiel 2024 du programme de seconde
générale (déposé dans ce Project) et tu t'y réfères systématiquement.

\medskip

Quand je te demande de générer un cours, tu structures ainsi~:
\begin{enumerate}[leftmargin=1.5em, label=\Roman*.]
  \item Objectifs notionnels (cite l'attendu de fin de programme correspondant).
  \item Pré-requis vérifiés.
  \item Activité d'introduction de 10 minutes.
  \item Cours formel structuré (définitions, propriétés, démonstrations adaptées).
  \item Trois exercices d'application gradués.
  \item Évaluation formative finale de 5 minutes.
\end{enumerate}

\medskip

Tu indiques toujours~: la durée totale estimée, les compétences travaillées
(modéliser, calculer, raisonner, communiquer), les écueils prévisibles pour
les élèves en difficulté.

\medskip

Ton ton est direct, pédagogique, sans jargon mathématique inutile. Tu
n'inventes jamais de référence à une page de manuel sans avoir vérifié
qu'elle existe dans les manuels déposés ici.
\end{skillbox}

\bordsec{L'usage au quotidien}

Le dimanche soir, Marc lance Claude Desktop, ouvre son Project, et écrit
simplement~:

\begin{quote}\itshape
« Génère un cours sur les fonctions affines. Durée 2~heures. Classe
hétérogène, sept élèves en grande difficulté. Cours du mardi à 8 heures —
ils seront fatigués. »
\end{quote}

Trente secondes plus tard, il a un plan de cours complet, structuré selon
le format qu'il a défini, ancré dans le Bulletin Officiel, avec des
exercices adaptés. Il relit, ajuste deux passages, ajoute son grain de sel
pédagogique sur l'introduction. Il a divisé par trois son temps de
préparation. La nuit récupérée, il l'utilise pour le suivi individuel
des élèves en difficulté — \textit{sa vraie valeur ajoutée}.

\pullquote{« L'IA me prépare le brouillon~; moi je fais la pédagogie.
C'est la première fois que la technologie me rend du temps au lieu de
m'en prendre. » \footnotesize\rmfamily\color{encrefaint} — Enseignant à Nouméa, retour anonyme~; secteur Éducation.}

\newpage

% -------------------------------------------------------------------------
% IV. HISTOIRE 2 — L'AVOCAT
% -------------------------------------------------------------------------

\dateline{IV. \quad Histoire 2}

\vspace{0.4cm}
{\color{encre}\rule{\linewidth}{0.4pt}}
\vspace{1cm}

\section*{L'avocat débordé}

\subsubsection*{Le contexte}

Maître Adam tient un cabinet généraliste à Nouméa~: contentieux civil,
droit de la famille, baux commerciaux. Une secrétaire deux jours par semaine,
le reste du temps c'est lui qui rédige, signe et expédie. Chaque dossier
demande~: une attestation à reformuler depuis les notes manuscrites du
client, un courrier de transmission, l'inventaire des pièces jointes
numérotées et indexées. Quatre à cinq heures par jour passées dans Word,
au lieu d'être avec ses clients.

\bordsec{Temps 1 — Son espace de travail}

Maître Adam crée un Project nommé \textit{« Cabinet — Rédaction »}~:

\begin{workspacebox}
\begin{itemize}[leftmargin=1.2em]
  \item Vingt \textbf{modèles de courriers} types du cabinet, sélectionnés
        comme étant \textit{ceux qui sonnent juste} (transmission de pièces,
        mise en demeure, demande de pièces, attestation, etc.).
  \item Un \textbf{glossaire interne} listant les formules de politesse,
        d'introduction, de conclusion qu'il utilise habituellement.
  \item Une \textbf{nomenclature des pièces} par type de dossier (en
        contentieux civil~: pièce~1 = exploit d'huissier, pièce~2 = courrier
        adverse, etc.).
  \item Le \textbf{Code civil} et le \textbf{Code de procédure civile} en
        version PDF officielle (Légifrance).
  \item Une note de cadrage~: \textit{« Ce que mon cabinet ne fait jamais »}
        — pour éviter les hallucinations de jurisprudence ou les renvois
        à des articles inexistants.
\end{itemize}
\end{workspacebox}

\bordsec{Temps 2 — Son skill}

\begin{skillbox}{Assistant de rédaction du cabinet}
Tu es l'assistant de rédaction du cabinet Adam à Nouméa. Ton rôle est
de produire des brouillons de documents juridiques (attestations, courriers,
inventaires de pièces) à partir des éléments bruts que je te donne — notes
manuscrites transcrites, audio dicté, e-mails du client.

\medskip

Pour chaque tâche~:

\begin{enumerate}[leftmargin=1.5em]
  \item Tu identifies le \textbf{type de document} attendu et tu reprends
        \textit{exactement} le style du modèle correspondant déposé dans
        ce Project.
  \item Tu reformules les éléments bruts dans le \textbf{registre juridique}
        approprié — sans jamais inventer de fait nouveau, ni d'article de
        loi, ni de date.
  \item En fin de document, tu listes les \textbf{pièces à joindre}
        numérotées selon notre nomenclature, et tu vérifies que chacune
        est référencée dans le corps du texte (indexation).
  \item Si une information manque (date, identité, montant, référence
        d'article), tu \textbf{poses la question avant de continuer}.
\end{enumerate}

\medskip

Tu ne cites jamais une jurisprudence ou un article du Code sans que je
te l'aie explicitement transmis. Tu n'utilises jamais d'emoji.
\end{skillbox}

\bordsec{L'usage au quotidien}

Lundi matin, Maître Adam sort de RDV avec une cliente. Elle a témoigné
oralement, il a pris des notes manuscrites pendant l'entretien. De retour
au cabinet, il dicte ses notes à voix haute dans son téléphone (cinq
minutes), récupère la transcription textuelle, la colle dans Claude
Desktop dans le Project \textit{Rédaction}, et écrit~:

\begin{quote}\itshape
« Voici le verbatim de l'entretien de ce matin avec Madame X. Rédige
l'attestation en bonne et due forme, dans le format des modèles déposés.
Liste les pièces à joindre selon notre nomenclature. »
\end{quote}

Une minute plus tard, il a~: une attestation propre, déjà mise en forme,
avec l'inventaire des pièces numérotées et chaque pièce référencée dans
le texte. Il relit, corrige une formulation, ajoute sa signature
électronique. La tâche qui prenait quarante minutes en prend dix.
Multiplié par trois ou quatre dossiers par jour, il récupère deux heures.
Deux heures qu'il consacre à ce qui ne peut pas être délégué — la
stratégie de défense, l'écoute du client, l'audience.

\pullquote{« Ce n'est pas un robot avocat. C'est un secrétariat ultra-rapide
et discipliné, qui me rend mes deux meilleures heures de la journée. »
\footnotesize\rmfamily\color{encrefaint} — Avocat à Nouméa, retour anonyme~;
secteur Profession libérale.}

\newpage

% -------------------------------------------------------------------------
% V. CONCLUSION
% -------------------------------------------------------------------------

\dateline{V. \quad Pour conclure}

\vspace{0.4cm}
{\color{encre}\rule{\linewidth}{0.4pt}}
\vspace{1cm}

\section*{Et pour vous~?}

Les deux histoires que vous venez de lire ne sont pas des cas exceptionnels.
Ce sont des configurations \textit{standards}, accessibles à n'importe quel
professionnel calédonien qui sait utiliser un ordinateur.

Le seul travail que demande Claude Desktop, c'est~:

\begin{itemize}
  \item \textbf{Une heure} pour bien construire l'espace de travail
        (choisir les bons documents, les déposer dans un Project).
  \item \textbf{Trente minutes} pour rédiger un skill propre, qui couvre
        votre métier et vos garde-fous.
  \item \textbf{Trois ou quatre semaines} de pratique quotidienne pour
        affiner — et sentir, vraiment, le temps revenir.
\end{itemize}

Ces trois investissements, vous pouvez les faire seul. Le guide est là
pour ça.

\bordsec{Si vous préférez le faire à deux}

Mon coaching individuel à Nouméa existe pour les professionnels qui veulent
\textit{compresser} ces trois étapes en une matinée. On regarde votre
métier, on construit votre espace de travail sur vos vrais dossiers, on
rédige votre skill ensemble. À midi, vous repartez avec quelque chose qui
tourne. Le mois suivant, je suis joignable par e-mail pour ajuster.

\vspace{0.7cm}

\subsubsection*{\textcolor{bordeaux}{Réservez un entretien découverte (gratuit, 30 minutes)}}

\vspace{0.3cm}

\begin{tabular}{@{}p{4cm}p{8cm}@{}}
\dateline{En ligne}     & \texttt{atelier-ia.ovh}\\[0.4em]
\dateline{Email}        & \texttt{tariel.vincent@gmail.com}\\[0.4em]
\dateline{Téléphone}    & \texttt{+687 95 07 86}\\[0.4em]
\dateline{Localisation} & Nouméa et Grand Nouméa, présentiel\\
\end{tabular}

\vfill

{\color{encre}\rule{\linewidth}{0.4pt}}
\vspace{0.3cm}

{\footnotesize\itshape\color{encrefaint}
Vincent Tariel — Docteur en mathématiques de l'École Polytechnique
(vision par ordinateur, 2009), agrégé de mathématiques, fondateur
de plusieurs startups en intelligence artificielle. Installé en
Nouvelle-Calédonie depuis 2020. Développe parallèlement la plateforme
pédagogique CommeUnJeu en utilisant Claude au quotidien — et transmet
ce qu'il apprend en coaching individuel.
}

\vspace{0.3cm}

\dateline{Fin du guide \quad\textbar\quad Atelier IA \quad\textbar\quad Édition 2026}

\end{document}
